% --- Archon physics formalization source begin ---
% archon:physics
% archon:covers PhyXMiniProblems/problem_phyx_mini_0918.lean
% archon:source-report reports/phyx_mini/problem_phyx_mini_0918.source.json

\chapter{Physics problem phyx\_mini\_0918}
\label{ch:PhyXMiniProblems_problem_phyx_mini_0918}

\paragraph{Problem source.}
\#\# Physical scenario

Figure shows a U-shaped conductor in a uniform magnetic field \$\textbackslash{}vec\{B\}\$ perpendicular to the plane of the figure and directed into the page. We lay a metal rod (the “slidewire”) with length \$L\$ across the two arms of the conductor, forming a circuit, and move it to the right with constant velocity \$\textbackslash{}vec\{v\}\$. This induces an emf and a current, which is why this device is called a slidewire generator.

\#\# Question

Find the magnitude and direction of the resulting induced emf.

\#\# Answer choices

- A: A: BLv

- B: B: \textbackslash{}frac\{Bv\}\{L\}

- C: C: -BLv

- D: D: \textbackslash{}frac\{L\}\{Bv\}

\#\# Figure caption (auxiliary; use the image as primary evidence)

The image depicts a physics diagram illustrating a rectangular loop of wire within a magnetic field. Here are the components and their relationships:

1. **Magnetic Field (\textbackslash{}(\textbackslash{}vec\{B\}\textbackslash{}))**: 
   - Represented by blue crosses, indicating the field direction into the plane of the screen.

2. **Rectangular Wire Loop**:
   - Consists of two vertical bars and two horizontal bars (one solid and one dashed).
   - The left side is fixed, while the right side is a movable bar.

3. **Current (\textbackslash{}(I\textbackslash{}))**:
   - Purple arrows on the sides of the loop indicate the direction of the induced current: counterclockwise at the top and clockwise at the bottom.

4. **Induced Electromotive Force (\textbackslash{}(\textbackslash{}mathcal\{E\}\textbackslash{}))**:
   - Indicated by a blue arrow pointing upwards on the left side, representing the direction of the force.

5. **Velocity (\textbackslash{}(v\textbackslash{}))**:
   - A green arrow on the right, movable bar, showing the direction of motion to the right.

6. **Displacement (\textbackslash{}(v \textbackslash{}, dt\textbackslash{}))**:
   - The horizontal black line above the loop shows the distance the bar moves in time \textbackslash{}(dt\textbackslash{}).

7. **Area Vector (\textbackslash{}(\textbackslash{}vec\{A\}\textbackslash{}))**:
   - A circle with an "X" inside represents the area vector perpendicular to the plane.

8. **Length (\textbackslash{}(L\textbackslash{}))**:
   - The vertical dimension of the loop is labeled with \textbackslash{}(L\textbackslash{}).

The diagram illustrates the concept of electromagnetic induction, where a change in magnetic flux through the loop induces an electromotive force and current due to the motion of the bar.

\#\# Dataset metadata

Electromagnetic Induction; Physical Model Grounding Reasoning, Multi-Formula Reasoning

\paragraph{Recorded answer/context.}
C: C: -BLv

\paragraph{Figure/image path.}
/root/proposal\_for\_physic/science-mango/phyx\_mini\_run/phyx\_data/test\_image/918.png

\paragraph{Formalization target.}
create a compiling Lean file with sorry bodies at `PhyXMiniProblems/problem\_phyx\_mini\_0918.lean`. The Lean declarations must preserve the physical quantities, dimensions or dimensional roles, figure labels, governing-law hypotheses, and final relation expressed by this problem.
Use Mathlib/Physlib names found through LeanExplore where available. If a physics API is missing, introduce faithful local abstractions rather than scalar placeholder aliases.

\begin{theorem}[Physics formalization target]
\label{thm:physics:phyx_mini_0918:target}
\lean{PhyXMiniProblems.ProblemPhyXMini0918.problem_phyx_mini_0918}
\uses{def:physics:phyx-mini-0918:phyxminiproblems-problemphyxmini0918-lengthinmeters, def:physics:phyx-mini-0918:phyxminiproblems-problemphyxmini0918-speedinmeterspersecond, def:physics:phyx-mini-0918:phyxminiproblems-problemphyxmini0918-magneticfluxdensityinteslas, def:physics:phyx-mini-0918:phyxminiproblems-problemphyxmini0918-signedemfinvolts, def:physics:phyx-mini-0918:phyxminiproblems-problemphyxmini0918-upwardunitvector, def:physics:phyx-mini-0918:phyxminiproblems-problemphyxmini0918-slidewiregeneratorsetup, def:physics:phyx-mini-0918:phyxminiproblems-problemphyxmini0918-matchesslidewiregeneratorscenario, def:physics:phyx-mini-0918:phyxminiproblems-problemphyxmini0918-matchessuppliedslidewirefigure, def:physics:phyx-mini-0918:phyxminiproblems-problemphyxmini0918-hasphysicalslidewireparameters, def:physics:phyx-mini-0918:phyxminiproblems-problemphyxmini0918-modelsuniformintopagemagneticfield, def:physics:phyx-mini-0918:phyxminiproblems-problemphyxmini0918-satisfiesslidewirekinematics, def:physics:phyx-mini-0918:phyxminiproblems-problemphyxmini0918-satisfiesuniformmagneticfluxlaw, def:physics:phyx-mini-0918:phyxminiproblems-problemphyxmini0918-satisfiesfaradayinductionlaw, def:physics:phyx-mini-0918:phyxminiproblems-problemphyxmini0918-satisfiesmotionalemfdirectionlaw, def:physics:phyx-mini-0918:phyxminiproblems-problemphyxmini0918-satisfiespassiveconductorresponse, def:physics:phyx-mini-0918:phyxminiproblems-problemphyxmini0918-recordeddatasetanswer}
The assigned autoformalize agent should translate this physics problem into Lean declarations in the covered file, with theorem and lemma proof bodies written as `by sorry`.
\end{theorem}
\begin{proof}
This is an autoformalization task, not a proof task. Produce statements that can later be proved by the physics prover without weakening the physical model.
\end{proof}

% --- Archon declaration topology begin ---
\section{Declaration topology}
\begin{definition}[magnetic Flux Density Dimension]
  \label{def:physics:phyx-mini-0918:phyxminiproblems-problemphyxmini0918-magneticfluxdensitydimension}
  \lean{PhyXMiniProblems.ProblemPhyXMini0918.magneticFluxDensityDimension}
  Magnetic flux density has dimension M T⁻¹ C⁻¹
\end{definition}
\begin{proof}
  This is the declared model definition of magnetic Flux Density Dimension.
\end{proof}
\begin{definition}[area Dimension]
  \label{def:physics:phyx-mini-0918:phyxminiproblems-problemphyxmini0918-areadimension}
  \lean{PhyXMiniProblems.ProblemPhyXMini0918.areaDimension}
  Area has dimension L²
\end{definition}
\begin{proof}
  This is the declared model definition of area Dimension.
\end{proof}
\begin{definition}[area Rate Dimension]
  \label{def:physics:phyx-mini-0918:phyxminiproblems-problemphyxmini0918-arearatedimension}
  \lean{PhyXMiniProblems.ProblemPhyXMini0918.areaRateDimension}
  \uses{def:physics:phyx-mini-0918:phyxminiproblems-problemphyxmini0918-areadimension}
  Area-change rate has dimension L² T⁻¹
\end{definition}
\begin{proof}
  This is the declared model definition of area Rate Dimension.
\end{proof}
\begin{definition}[magnetic Flux Dimension]
  \label{def:physics:phyx-mini-0918:phyxminiproblems-problemphyxmini0918-magneticfluxdimension}
  \lean{PhyXMiniProblems.ProblemPhyXMini0918.magneticFluxDimension}
  \uses{def:physics:phyx-mini-0918:phyxminiproblems-problemphyxmini0918-magneticfluxdensitydimension, def:physics:phyx-mini-0918:phyxminiproblems-problemphyxmini0918-areadimension}
  Magnetic flux has dimension M L² T⁻¹ C⁻¹
\end{definition}
\begin{proof}
  This is the declared model definition of magnetic Flux Dimension.
\end{proof}
\begin{definition}[electromotive Force Dimension]
  \label{def:physics:phyx-mini-0918:phyxminiproblems-problemphyxmini0918-electromotiveforcedimension}
  \lean{PhyXMiniProblems.ProblemPhyXMini0918.electromotiveForceDimension}
  \uses{def:physics:phyx-mini-0918:phyxminiproblems-problemphyxmini0918-magneticfluxdimension}
  Electromotive force, or magnetic-flux rate, has dimension M L² T⁻² C⁻¹
\end{definition}
\begin{proof}
  This is the declared model definition of electromotive Force Dimension.
\end{proof}
\begin{definition}[electric Current Dimension]
  \label{def:physics:phyx-mini-0918:phyxminiproblems-problemphyxmini0918-electriccurrentdimension}
  \lean{PhyXMiniProblems.ProblemPhyXMini0918.electricCurrentDimension}
  Electric current has dimension C T⁻¹
\end{definition}
\begin{proof}
  This is the declared model definition of electric Current Dimension.
\end{proof}
\begin{definition}[Length Quantity]
  \label{def:physics:phyx-mini-0918:phyxminiproblems-problemphyxmini0918-lengthquantity}
  \lean{PhyXMiniProblems.ProblemPhyXMini0918.LengthQuantity}
  A nonnegative, unit-independent physical length
\end{definition}
\begin{proof}
  This is the declared model definition of Length Quantity.
\end{proof}
\begin{definition}[Time Quantity]
  \label{def:physics:phyx-mini-0918:phyxminiproblems-problemphyxmini0918-timequantity}
  \lean{PhyXMiniProblems.ProblemPhyXMini0918.TimeQuantity}
  A nonnegative, unit-independent elapsed time
\end{definition}
\begin{proof}
  This is the declared model definition of Time Quantity.
\end{proof}
\begin{definition}[Speed Quantity]
  \label{def:physics:phyx-mini-0918:phyxminiproblems-problemphyxmini0918-speedquantity}
  \lean{PhyXMiniProblems.ProblemPhyXMini0918.SpeedQuantity}
  A nonnegative, unit-independent speed
\end{definition}
\begin{proof}
  This is the declared model definition of Speed Quantity.
\end{proof}
\begin{definition}[Area Quantity]
  \label{def:physics:phyx-mini-0918:phyxminiproblems-problemphyxmini0918-areaquantity}
  \lean{PhyXMiniProblems.ProblemPhyXMini0918.AreaQuantity}
  \uses{def:physics:phyx-mini-0918:phyxminiproblems-problemphyxmini0918-areadimension}
  A nonnegative, unit-independent area
\end{definition}
\begin{proof}
  This is the declared model definition of Area Quantity.
\end{proof}
\begin{definition}[Signed Area Rate Quantity]
  \label{def:physics:phyx-mini-0918:phyxminiproblems-problemphyxmini0918-signedarearatequantity}
  \lean{PhyXMiniProblems.ProblemPhyXMini0918.SignedAreaRateQuantity}
  \uses{def:physics:phyx-mini-0918:phyxminiproblems-problemphyxmini0918-arearatedimension}
  A signed area-change rate relative to the chosen area orientation
\end{definition}
\begin{proof}
  This is the declared model definition of Signed Area Rate Quantity.
\end{proof}
\begin{definition}[Magnetic Flux Density Quantity]
  \label{def:physics:phyx-mini-0918:phyxminiproblems-problemphyxmini0918-magneticfluxdensityquantity}
  \lean{PhyXMiniProblems.ProblemPhyXMini0918.MagneticFluxDensityQuantity}
  \uses{def:physics:phyx-mini-0918:phyxminiproblems-problemphyxmini0918-magneticfluxdensitydimension}
  A nonnegative magnetic-flux-density magnitude
\end{definition}
\begin{proof}
  This is the declared model definition of Magnetic Flux Density Quantity.
\end{proof}
\begin{definition}[Signed Magnetic Flux Quantity]
  \label{def:physics:phyx-mini-0918:phyxminiproblems-problemphyxmini0918-signedmagneticfluxquantity}
  \lean{PhyXMiniProblems.ProblemPhyXMini0918.SignedMagneticFluxQuantity}
  \uses{def:physics:phyx-mini-0918:phyxminiproblems-problemphyxmini0918-magneticfluxdimension}
  Signed magnetic flux relative to the chosen area normal
\end{definition}
\begin{proof}
  This is the declared model definition of Signed Magnetic Flux Quantity.
\end{proof}
\begin{definition}[Signed Magnetic Flux Rate Quantity]
  \label{def:physics:phyx-mini-0918:phyxminiproblems-problemphyxmini0918-signedmagneticfluxratequantity}
  \lean{PhyXMiniProblems.ProblemPhyXMini0918.SignedMagneticFluxRateQuantity}
  \uses{def:physics:phyx-mini-0918:phyxminiproblems-problemphyxmini0918-electromotiveforcedimension}
  Signed magnetic-flux change rate
\end{definition}
\begin{proof}
  This is the declared model definition of Signed Magnetic Flux Rate Quantity.
\end{proof}
\begin{definition}[Signed Emf Quantity]
  \label{def:physics:phyx-mini-0918:phyxminiproblems-problemphyxmini0918-signedemfquantity}
  \lean{PhyXMiniProblems.ProblemPhyXMini0918.SignedEmfQuantity}
  \uses{def:physics:phyx-mini-0918:phyxminiproblems-problemphyxmini0918-electromotiveforcedimension}
  Signed electromotive force relative to the positive loop orientation
\end{definition}
\begin{proof}
  This is the declared model definition of Signed Emf Quantity.
\end{proof}
\begin{definition}[Current Magnitude Quantity]
  \label{def:physics:phyx-mini-0918:phyxminiproblems-problemphyxmini0918-currentmagnitudequantity}
  \lean{PhyXMiniProblems.ProblemPhyXMini0918.CurrentMagnitudeQuantity}
  \uses{def:physics:phyx-mini-0918:phyxminiproblems-problemphyxmini0918-electriccurrentdimension}
  A nonnegative electric-current magnitude
\end{definition}
\begin{proof}
  This is the declared model definition of Current Magnitude Quantity.
\end{proof}
\begin{definition}[length In Meters]
  \label{def:physics:phyx-mini-0918:phyxminiproblems-problemphyxmini0918-lengthinmeters}
  \lean{PhyXMiniProblems.ProblemPhyXMini0918.lengthInMeters}
  \uses{def:physics:phyx-mini-0918:phyxminiproblems-problemphyxmini0918-lengthquantity}
  Read a length in coherent-SI metres
\end{definition}
\begin{proof}
  This is the declared model definition of length In Meters.
\end{proof}
\begin{definition}[time In Seconds]
  \label{def:physics:phyx-mini-0918:phyxminiproblems-problemphyxmini0918-timeinseconds}
  \lean{PhyXMiniProblems.ProblemPhyXMini0918.timeInSeconds}
  \uses{def:physics:phyx-mini-0918:phyxminiproblems-problemphyxmini0918-timequantity}
  Read elapsed time in coherent-SI seconds
\end{definition}
\begin{proof}
  This is the declared model definition of time In Seconds.
\end{proof}
\begin{definition}[speed In Meters Per Second]
  \label{def:physics:phyx-mini-0918:phyxminiproblems-problemphyxmini0918-speedinmeterspersecond}
  \lean{PhyXMiniProblems.ProblemPhyXMini0918.speedInMetersPerSecond}
  \uses{def:physics:phyx-mini-0918:phyxminiproblems-problemphyxmini0918-speedquantity}
  Read speed in coherent-SI metres per second
\end{definition}
\begin{proof}
  This is the declared model definition of speed In Meters Per Second.
\end{proof}
\begin{definition}[area In Square Meters]
  \label{def:physics:phyx-mini-0918:phyxminiproblems-problemphyxmini0918-areainsquaremeters}
  \lean{PhyXMiniProblems.ProblemPhyXMini0918.areaInSquareMeters}
  \uses{def:physics:phyx-mini-0918:phyxminiproblems-problemphyxmini0918-areaquantity}
  Read area in coherent-SI square metres
\end{definition}
\begin{proof}
  This is the declared model definition of area In Square Meters.
\end{proof}
\begin{definition}[area Rate In Square Meters Per Second]
  \label{def:physics:phyx-mini-0918:phyxminiproblems-problemphyxmini0918-arearateinsquaremeterspersecond}
  \lean{PhyXMiniProblems.ProblemPhyXMini0918.areaRateInSquareMetersPerSecond}
  \uses{def:physics:phyx-mini-0918:phyxminiproblems-problemphyxmini0918-signedarearatequantity}
  Read signed area-change rate in square metres per second
\end{definition}
\begin{proof}
  This is the declared model definition of area Rate In Square Meters Per Second.
\end{proof}
\begin{definition}[magnetic Flux Density In Teslas]
  \label{def:physics:phyx-mini-0918:phyxminiproblems-problemphyxmini0918-magneticfluxdensityinteslas}
  \lean{PhyXMiniProblems.ProblemPhyXMini0918.magneticFluxDensityInTeslas}
  \uses{def:physics:phyx-mini-0918:phyxminiproblems-problemphyxmini0918-magneticfluxdensityquantity}
  Read magnetic-flux-density magnitude in coherent-SI teslas
\end{definition}
\begin{proof}
  This is the declared model definition of magnetic Flux Density In Teslas.
\end{proof}
\begin{definition}[magnetic Flux In Webers]
  \label{def:physics:phyx-mini-0918:phyxminiproblems-problemphyxmini0918-magneticfluxinwebers}
  \lean{PhyXMiniProblems.ProblemPhyXMini0918.magneticFluxInWebers}
  \uses{def:physics:phyx-mini-0918:phyxminiproblems-problemphyxmini0918-signedmagneticfluxquantity}
  Read signed magnetic flux in coherent-SI webers
\end{definition}
\begin{proof}
  This is the declared model definition of magnetic Flux In Webers.
\end{proof}
\begin{definition}[magnetic Flux Rate In Webers Per Second]
  \label{def:physics:phyx-mini-0918:phyxminiproblems-problemphyxmini0918-magneticfluxrateinweberspersecond}
  \lean{PhyXMiniProblems.ProblemPhyXMini0918.magneticFluxRateInWebersPerSecond}
  \uses{def:physics:phyx-mini-0918:phyxminiproblems-problemphyxmini0918-signedmagneticfluxratequantity}
  Read signed magnetic-flux rate in coherent-SI webers per second
\end{definition}
\begin{proof}
  This is the declared model definition of magnetic Flux Rate In Webers Per Second.
\end{proof}
\begin{definition}[signed Emf In Volts]
  \label{def:physics:phyx-mini-0918:phyxminiproblems-problemphyxmini0918-signedemfinvolts}
  \lean{PhyXMiniProblems.ProblemPhyXMini0918.signedEmfInVolts}
  \uses{def:physics:phyx-mini-0918:phyxminiproblems-problemphyxmini0918-signedemfquantity}
  Read signed electromotive force in coherent-SI volts
\end{definition}
\begin{proof}
  This is the declared model definition of signed Emf In Volts.
\end{proof}
\begin{definition}[current Magnitude In Amperes]
  \label{def:physics:phyx-mini-0918:phyxminiproblems-problemphyxmini0918-currentmagnitudeinamperes}
  \lean{PhyXMiniProblems.ProblemPhyXMini0918.currentMagnitudeInAmperes}
  \uses{def:physics:phyx-mini-0918:phyxminiproblems-problemphyxmini0918-currentmagnitudequantity}
  Read current magnitude in coherent-SI amperes
\end{definition}
\begin{proof}
  This is the declared model definition of current Magnitude In Amperes.
\end{proof}
\begin{definition}[Direction Vector]
  \label{def:physics:phyx-mini-0918:phyxminiproblems-problemphyxmini0918-directionvector}
  \lean{PhyXMiniProblems.ProblemPhyXMini0918.DirectionVector}
  A dimensionless direction vector in physical three-space
\end{definition}
\begin{proof}
  This is the declared model definition of Direction Vector.
\end{proof}
\begin{definition}[rightward Unit Vector]
  \label{def:physics:phyx-mini-0918:phyxminiproblems-problemphyxmini0918-rightwardunitvector}
  \lean{PhyXMiniProblems.ProblemPhyXMini0918.rightwardUnitVector}
  \uses{def:physics:phyx-mini-0918:phyxminiproblems-problemphyxmini0918-directionvector}
  Unit direction pointing right in the plane of the diagram
\end{definition}
\begin{proof}
  This is the declared model definition of rightward Unit Vector.
\end{proof}
\begin{definition}[upward Unit Vector]
  \label{def:physics:phyx-mini-0918:phyxminiproblems-problemphyxmini0918-upwardunitvector}
  \lean{PhyXMiniProblems.ProblemPhyXMini0918.upwardUnitVector}
  \uses{def:physics:phyx-mini-0918:phyxminiproblems-problemphyxmini0918-directionvector}
  Unit direction pointing upward in the plane of the diagram
\end{definition}
\begin{proof}
  This is the declared model definition of upward Unit Vector.
\end{proof}
\begin{definition}[into Page Unit Vector]
  \label{def:physics:phyx-mini-0918:phyxminiproblems-problemphyxmini0918-intopageunitvector}
  \lean{PhyXMiniProblems.ProblemPhyXMini0918.intoPageUnitVector}
  \uses{def:physics:phyx-mini-0918:phyxminiproblems-problemphyxmini0918-directionvector}
  Unit direction pointing into the page
\end{definition}
\begin{proof}
  This is the declared model definition of into Page Unit Vector.
\end{proof}
\begin{definition}[euclidean Cross Product]
  \label{def:physics:phyx-mini-0918:phyxminiproblems-problemphyxmini0918-euclideancrossproduct}
  \lean{PhyXMiniProblems.ProblemPhyXMini0918.euclideanCrossProduct}
  \uses{def:physics:phyx-mini-0918:phyxminiproblems-problemphyxmini0918-directionvector}
  The ordinary right-handed cross product on Euclidean three-vectors. This wrapper transports Mathlib's crossProduct through the EuclideanSpace representation used by Physlib
\end{definition}
\begin{proof}
  This is the declared model definition of euclidean Cross Product.
\end{proof}
\begin{definition}[Loop Direction]
  \label{def:physics:phyx-mini-0918:phyxminiproblems-problemphyxmini0918-loopdirection}
  \lean{PhyXMiniProblems.ProblemPhyXMini0918.LoopDirection}
  The two orientations around the rectangular conducting loop
\end{definition}
\begin{proof}
  This is the declared model definition of Loop Direction.
\end{proof}
\begin{definition}[slidewire Direction Vector]
  \label{def:physics:phyx-mini-0918:phyxminiproblems-problemphyxmini0918-slidewiredirectionvector}
  \lean{PhyXMiniProblems.ProblemPhyXMini0918.slidewireDirectionVector}
  \uses{def:physics:phyx-mini-0918:phyxminiproblems-problemphyxmini0918-directionvector, def:physics:phyx-mini-0918:phyxminiproblems-problemphyxmini0918-upwardunitvector, def:physics:phyx-mini-0918:phyxminiproblems-problemphyxmini0918-loopdirection}
  Direction along the moving right-hand rod for each loop orientation
\end{definition}
\begin{proof}
  This is the declared model definition of slidewire Direction Vector.
\end{proof}
\begin{definition}[Rail]
  \label{def:physics:phyx-mini-0918:phyxminiproblems-problemphyxmini0918-rail}
  \lean{PhyXMiniProblems.ProblemPhyXMini0918.Rail}
  The two horizontal rails, distinguished by vertical placement
\end{definition}
\begin{proof}
  This is the declared model definition of Rail.
\end{proof}
\begin{definition}[Diagram Direction]
  \label{def:physics:phyx-mini-0918:phyxminiproblems-problemphyxmini0918-diagramdirection}
  \lean{PhyXMiniProblems.ProblemPhyXMini0918.DiagramDirection}
  Planar and normal directions used to transcribe the raster
\end{definition}
\begin{proof}
  This is the declared model definition of Diagram Direction.
\end{proof}
\begin{definition}[Figure Object]
  \label{def:physics:phyx-mini-0918:phyxminiproblems-problemphyxmini0918-figureobject}
  \lean{PhyXMiniProblems.ProblemPhyXMini0918.FigureObject}
  Physical and graphical objects visible in 918.png
\end{definition}
\begin{proof}
  This is the declared model definition of Figure Object.
\end{proof}
\begin{definition}[Figure Label]
  \label{def:physics:phyx-mini-0918:phyxminiproblems-problemphyxmini0918-figurelabel}
  \lean{PhyXMiniProblems.ProblemPhyXMini0918.FigureLabel}
  Symbolic labels printed in 918.png
\end{definition}
\begin{proof}
  This is the declared model definition of Figure Label.
\end{proof}
\begin{definition}[Normal Vector Glyph]
  \label{def:physics:phyx-mini-0918:phyxminiproblems-problemphyxmini0918-normalvectorglyph}
  \lean{PhyXMiniProblems.ProblemPhyXMini0918.NormalVectorGlyph}
  Glyph used for a vector normal to the page
\end{definition}
\begin{proof}
  This is the declared model definition of Normal Vector Glyph.
\end{proof}
\begin{definition}[Slidewire Generator Figure]
  \label{def:physics:phyx-mini-0918:phyxminiproblems-problemphyxmini0918-slidewiregeneratorfigure}
  \lean{PhyXMiniProblems.ProblemPhyXMini0918.SlidewireGeneratorFigure}
  \uses{def:physics:phyx-mini-0918:phyxminiproblems-problemphyxmini0918-rail, def:physics:phyx-mini-0918:phyxminiproblems-problemphyxmini0918-diagramdirection, def:physics:phyx-mini-0918:phyxminiproblems-problemphyxmini0918-figureobject, def:physics:phyx-mini-0918:phyxminiproblems-problemphyxmini0918-figurelabel, def:physics:phyx-mini-0918:phyxminiproblems-problemphyxmini0918-normalvectorglyph}
  Typed presentation data from the primary raster. The emf and current arrows are recorded as visible, but their answer-bearing directions are deliberately not stored here; those directions remain conclusions of the governing laws
\end{definition}
\begin{proof}
  This is the declared model definition of Slidewire Generator Figure.
\end{proof}
\begin{definition}[Conductor Model]
  \label{def:physics:phyx-mini-0918:phyxminiproblems-problemphyxmini0918-conductormodel}
  \lean{PhyXMiniProblems.ProblemPhyXMini0918.ConductorModel}
  Material idealization of the rails and slidewire
\end{definition}
\begin{proof}
  This is the declared model definition of Conductor Model.
\end{proof}
\begin{definition}[Slidewire Geometry Model]
  \label{def:physics:phyx-mini-0918:phyxminiproblems-problemphyxmini0918-slidewiregeometrymodel}
  \lean{PhyXMiniProblems.ProblemPhyXMini0918.SlidewireGeometryModel}
  Geometric idealization of the moving bar and rails
\end{definition}
\begin{proof}
  This is the declared model definition of Slidewire Geometry Model.
\end{proof}
\begin{definition}[Slidewire Generator Setup]
  \label{def:physics:phyx-mini-0918:phyxminiproblems-problemphyxmini0918-slidewiregeneratorsetup}
  \lean{PhyXMiniProblems.ProblemPhyXMini0918.SlidewireGeneratorSetup}
  \uses{def:physics:phyx-mini-0918:phyxminiproblems-problemphyxmini0918-lengthquantity, def:physics:phyx-mini-0918:phyxminiproblems-problemphyxmini0918-timequantity, def:physics:phyx-mini-0918:phyxminiproblems-problemphyxmini0918-speedquantity, def:physics:phyx-mini-0918:phyxminiproblems-problemphyxmini0918-areaquantity, def:physics:phyx-mini-0918:phyxminiproblems-problemphyxmini0918-signedarearatequantity, def:physics:phyx-mini-0918:phyxminiproblems-problemphyxmini0918-magneticfluxdensityquantity, def:physics:phyx-mini-0918:phyxminiproblems-problemphyxmini0918-signedmagneticfluxquantity, def:physics:phyx-mini-0918:phyxminiproblems-problemphyxmini0918-signedmagneticfluxratequantity, def:physics:phyx-mini-0918:phyxminiproblems-problemphyxmini0918-signedemfquantity, def:physics:phyx-mini-0918:phyxminiproblems-problemphyxmini0918-currentmagnitudequantity, def:physics:phyx-mini-0918:phyxminiproblems-problemphyxmini0918-directionvector, def:physics:phyx-mini-0918:phyxminiproblems-problemphyxmini0918-loopdirection, def:physics:phyx-mini-0918:phyxminiproblems-problemphyxmini0918-slidewiregeneratorfigure, def:physics:phyx-mini-0918:phyxminiproblems-problemphyxmini0918-conductormodel, def:physics:phyx-mini-0918:phyxminiproblems-problemphyxmini0918-slidewiregeometrymodel}
  The apparatus and its independent observables. The emf, flux rate, area rate, and directions are not defined from B, L, v, or any answer choice
\end{definition}
\begin{proof}
  This is the declared model definition of Slidewire Generator Setup.
\end{proof}
\begin{definition}[Matches Slidewire Generator Scenario]
  \label{def:physics:phyx-mini-0918:phyxminiproblems-problemphyxmini0918-matchesslidewiregeneratorscenario}
  \lean{PhyXMiniProblems.ProblemPhyXMini0918.MatchesSlidewireGeneratorScenario}
  \uses{def:physics:phyx-mini-0918:phyxminiproblems-problemphyxmini0918-rightwardunitvector, def:physics:phyx-mini-0918:phyxminiproblems-problemphyxmini0918-intopageunitvector, def:physics:phyx-mini-0918:phyxminiproblems-problemphyxmini0918-slidewiregeneratorsetup}
  Prose-level apparatus data, excluding the requested emf result
\end{definition}
\begin{proof}
  This is the declared model definition of Matches Slidewire Generator Scenario.
\end{proof}
\begin{definition}[Matches Supplied Slidewire Figure]
  \label{def:physics:phyx-mini-0918:phyxminiproblems-problemphyxmini0918-matchessuppliedslidewirefigure}
  \lean{PhyXMiniProblems.ProblemPhyXMini0918.MatchesSuppliedSlidewireFigure}
  \uses{def:physics:phyx-mini-0918:phyxminiproblems-problemphyxmini0918-slidewiregeneratorsetup}
  Literal geometry and non-answer annotations visible in 918.png
\end{definition}
\begin{proof}
  This is the declared model definition of Matches Supplied Slidewire Figure.
\end{proof}
\begin{definition}[Has Physical Slidewire Parameters]
  \label{def:physics:phyx-mini-0918:phyxminiproblems-problemphyxmini0918-hasphysicalslidewireparameters}
  \lean{PhyXMiniProblems.ProblemPhyXMini0918.HasPhysicalSlidewireParameters}
  \uses{def:physics:phyx-mini-0918:phyxminiproblems-problemphyxmini0918-lengthinmeters, def:physics:phyx-mini-0918:phyxminiproblems-problemphyxmini0918-speedinmeterspersecond, def:physics:phyx-mini-0918:phyxminiproblems-problemphyxmini0918-magneticfluxdensityinteslas, def:physics:phyx-mini-0918:phyxminiproblems-problemphyxmini0918-slidewiregeneratorsetup}
  Positivity and nondegeneracy of the physical parameters in B L v
\end{definition}
\begin{proof}
  This is the declared model definition of Has Physical Slidewire Parameters.
\end{proof}
\begin{definition}[Models Uniform Into Page Magnetic Field]
  \label{def:physics:phyx-mini-0918:phyxminiproblems-problemphyxmini0918-modelsuniformintopagemagneticfield}
  \lean{PhyXMiniProblems.ProblemPhyXMini0918.ModelsUniformIntoPageMagneticField}
  \uses{def:physics:phyx-mini-0918:phyxminiproblems-problemphyxmini0918-magneticfluxdensityinteslas, def:physics:phyx-mini-0918:phyxminiproblems-problemphyxmini0918-intopageunitvector, def:physics:phyx-mini-0918:phyxminiproblems-problemphyxmini0918-slidewiregeneratorsetup}
  The full Physlib magnetic field is spatially and temporally uniform, directed into the page, and has the given magnetic-flux-density readout
\end{definition}
\begin{proof}
  This is the declared model definition of Models Uniform Into Page Magnetic Field.
\end{proof}
\begin{definition}[Satisfies Slidewire Kinematics]
  \label{def:physics:phyx-mini-0918:phyxminiproblems-problemphyxmini0918-satisfiesslidewirekinematics}
  \lean{PhyXMiniProblems.ProblemPhyXMini0918.SatisfiesSlidewireKinematics}
  \uses{def:physics:phyx-mini-0918:phyxminiproblems-problemphyxmini0918-timequantity, def:physics:phyx-mini-0918:phyxminiproblems-problemphyxmini0918-lengthinmeters, def:physics:phyx-mini-0918:phyxminiproblems-problemphyxmini0918-timeinseconds, def:physics:phyx-mini-0918:phyxminiproblems-problemphyxmini0918-speedinmeterspersecond, def:physics:phyx-mini-0918:phyxminiproblems-problemphyxmini0918-arearateinsquaremeterspersecond, def:physics:phyx-mini-0918:phyxminiproblems-problemphyxmini0918-slidewiregeneratorsetup}
  Constant-speed kinematics: the displacement label means v dt, and the rectangular loop area grows at rate L v. Neither statement mentions flux or emf
\end{definition}
\begin{proof}
  This is the declared model definition of Satisfies Slidewire Kinematics.
\end{proof}
\begin{definition}[Satisfies Uniform Magnetic Flux Law]
  \label{def:physics:phyx-mini-0918:phyxminiproblems-problemphyxmini0918-satisfiesuniformmagneticfluxlaw}
  \lean{PhyXMiniProblems.ProblemPhyXMini0918.SatisfiesUniformMagneticFluxLaw}
  \uses{def:physics:phyx-mini-0918:phyxminiproblems-problemphyxmini0918-areainsquaremeters, def:physics:phyx-mini-0918:phyxminiproblems-problemphyxmini0918-arearateinsquaremeterspersecond, def:physics:phyx-mini-0918:phyxminiproblems-problemphyxmini0918-magneticfluxdensityinteslas, def:physics:phyx-mini-0918:phyxminiproblems-problemphyxmini0918-magneticfluxinwebers, def:physics:phyx-mini-0918:phyxminiproblems-problemphyxmini0918-magneticfluxrateinweberspersecond, def:physics:phyx-mini-0918:phyxminiproblems-problemphyxmini0918-slidewiregeneratorsetup}
  For a uniform field parallel to the chosen area normal, flux is B A and its rate is B dA/dt. The law does not mention emf or substitute the slidewire kinematics into the flux rate
\end{definition}
\begin{proof}
  This is the declared model definition of Satisfies Uniform Magnetic Flux Law.
\end{proof}
\begin{definition}[Satisfies Faraday Induction Law]
  \label{def:physics:phyx-mini-0918:phyxminiproblems-problemphyxmini0918-satisfiesfaradayinductionlaw}
  \lean{PhyXMiniProblems.ProblemPhyXMini0918.SatisfiesFaradayInductionLaw}
  \uses{def:physics:phyx-mini-0918:phyxminiproblems-problemphyxmini0918-magneticfluxrateinweberspersecond, def:physics:phyx-mini-0918:phyxminiproblems-problemphyxmini0918-signedemfinvolts, def:physics:phyx-mini-0918:phyxminiproblems-problemphyxmini0918-slidewiregeneratorsetup}
  Faraday's law relative to the selected positive area normal and boundary orientation. It contains neither the rod length nor its speed
\end{definition}
\begin{proof}
  This is the declared model definition of Satisfies Faraday Induction Law.
\end{proof}
\begin{definition}[Satisfies Motional Emf Direction Law]
  \label{def:physics:phyx-mini-0918:phyxminiproblems-problemphyxmini0918-satisfiesmotionalemfdirectionlaw}
  \lean{PhyXMiniProblems.ProblemPhyXMini0918.SatisfiesMotionalEmfDirectionLaw}
  \uses{def:physics:phyx-mini-0918:phyxminiproblems-problemphyxmini0918-euclideancrossproduct, def:physics:phyx-mini-0918:phyxminiproblems-problemphyxmini0918-slidewiredirectionvector, def:physics:phyx-mini-0918:phyxminiproblems-problemphyxmini0918-slidewiregeneratorsetup}
  The magnetic Lorentz-force direction is v × B. Loop geometry states how each possible orientation traverses the rod, but does not select an answer
\end{definition}
\begin{proof}
  This is the declared model definition of Satisfies Motional Emf Direction Law.
\end{proof}
\begin{definition}[Satisfies Passive Conductor Response]
  \label{def:physics:phyx-mini-0918:phyxminiproblems-problemphyxmini0918-satisfiespassiveconductorresponse}
  \lean{PhyXMiniProblems.ProblemPhyXMini0918.SatisfiesPassiveConductorResponse}
  \uses{def:physics:phyx-mini-0918:phyxminiproblems-problemphyxmini0918-slidewiregeneratorsetup}
  A passive conducting loop carries conventional current in the same loop orientation as its induced emf. Its magnitude remains unspecified because the resistance is not supplied
\end{definition}
\begin{proof}
  This is the declared model definition of Satisfies Passive Conductor Response.
\end{proof}
\begin{definition}[Answer Choice]
  \label{def:physics:phyx-mini-0918:phyxminiproblems-problemphyxmini0918-answerchoice}
  \lean{PhyXMiniProblems.ProblemPhyXMini0918.AnswerChoice}
  The four answer labels printed with the question
\end{definition}
\begin{proof}
  This is the declared model definition of Answer Choice.
\end{proof}
\begin{definition}[displayed Signed Emf In Volts]
  \label{def:physics:phyx-mini-0918:phyxminiproblems-problemphyxmini0918-answerchoice-displayedsignedemfinvolts}
  \lean{PhyXMiniProblems.ProblemPhyXMini0918.AnswerChoice.displayedSignedEmfInVolts}
  \uses{def:physics:phyx-mini-0918:phyxminiproblems-problemphyxmini0918-lengthinmeters, def:physics:phyx-mini-0918:phyxminiproblems-problemphyxmini0918-speedinmeterspersecond, def:physics:phyx-mini-0918:phyxminiproblems-problemphyxmini0918-magneticfluxdensityinteslas, def:physics:phyx-mini-0918:phyxminiproblems-problemphyxmini0918-slidewiregeneratorsetup, def:physics:phyx-mini-0918:phyxminiproblems-problemphyxmini0918-answerchoice}
  The signed-emf expression, in volts, printed for each answer choice
\end{definition}
\begin{proof}
  This is the declared model definition of displayed Signed Emf In Volts.
\end{proof}
\begin{definition}[recorded Dataset Answer]
  \label{def:physics:phyx-mini-0918:phyxminiproblems-problemphyxmini0918-recordeddatasetanswer}
  \lean{PhyXMiniProblems.ProblemPhyXMini0918.recordedDatasetAnswer}
  \uses{def:physics:phyx-mini-0918:phyxminiproblems-problemphyxmini0918-answerchoice}
  The answer label recorded by the source dataset
\end{definition}
\begin{proof}
  This is the declared model definition of recorded Dataset Answer.
\end{proof}
% --- Archon declaration topology end ---
% --- Archon physics formalization source end ---
